\documentclass[uplatex,dvipdfmx,11pt,a4paper]{jsarticle}

\usepackage[a4paper,margin=24mm]{geometry}
\usepackage[dvipdfmx]{graphicx}
\usepackage{booktabs}
\usepackage{longtable}
\usepackage{amsmath}
\usepackage{caption}
\usepackage{float}
\usepackage{xcolor}
\usepackage{url}
\usepackage[dvipdfmx]{hyperref}
\usepackage{pxjahyper}

\hypersetup{
  colorlinks=true,
  linkcolor=blue,
  urlcolor=blue,
  citecolor=blue,
  pdftitle={法然・親鸞・浄土教祖師文献の意味地図：三層分析に向けたSAT漢文ベースの予備的検討},
  pdfauthor={上山大信},
  pdfkeywords={法然, 親鸞, 選択本願念仏集, 教行信証, 浄土教, 意味埋め込み, 三層分析, SAT, text-embedding-3-large}
}

\captionsetup{font=small,labelfont=bf}
\renewcommand{\figurename}{図}
\renewcommand{\tablename}{表}
\renewcommand{\refname}{参考文献}
\renewcommand{\abstractname}{要旨}

\title{法然・親鸞・浄土教祖師文献の意味地図：\\[0.35em]三層分析に向けたSAT漢文ベースの予備的検討}
\author{上山大信}
\date{2026年6月4日草稿}

\begin{document}

\maketitle

\begin{abstract}
本稿は、法然『選擇本願念佛集』と親鸞『顯淨土眞實教行證文類』を中心に、浄土教祖師文献を重ねた意味埋め込み地図を作成し、三層分析の基礎を整えるための予備的報告である。筆者は先に「意味・文体・典拠マーカーの三層地図による仏教文献探索：宗派別参照群・阿弥陀経二訳・親鸞文献の予備的分析」において、宗派別参照テキスト群、阿弥陀経二訳、親鸞文献を対象に三層地図を試みた\cite{ueyama2026three-layer-map}。本稿ではその課題を法然・親鸞・七高僧系文献へ絞り、SAT漢文本文、Unicode-safeな700トークンチャンク、100トークン重なり、OpenAI \texttt{text-embedding-3-large} による埋め込みを用いた。

結果として、法然チャンク69件の最近傍非自己は親鸞50件、道綽11件、源信4件、善導3件、曇鸞1件であり、法然と親鸞の強い重なりが確認された。一方、親鸞チャンク191件の最近傍非自己は法然48件、道綽45件、源信36件、曇鸞23件、天親21件、善導17件、龍樹1件に分かれた。これは、親鸞が法然と重なりつつ、道綽・源信・曇鸞・天親など複数の祖師文献方向へ広がることを示す。さらに法然のはみ出し領域は、念仏の意味そのものの差というより、正雑二行、諸行との取捨、本願、付属・証誠などを用いて念仏を選び出す論証のまとまりとして現れた。これに対し親鸞のはみ出し領域は、信巻の信・三心・罪救済、化身土巻の護法・鬼神・宇宙秩序、外教批判、方便・真仮整理に集中した。

本稿の作業仮説は、法然は「念仏を選ぶ論理」を構成し、親鸞はその念仏を「本願・信・証・真仏土/化身土の典拠的体系」として展開する、というものである。ただし、この仮説は埋め込み、近傍関係、語群カウント、典拠マーカーから得た探索的推測であり、歴史的研究成果を参照して裏づけたものではない。引用文の立場と著者自身の立場を分離する作業、巻別・章別境界の精密化、古写経・異本研究との照合は今後の課題である。
\end{abstract}

\noindent\textbf{キーワード：} 法然，親鸞，選択本願念仏集，教行信証，浄土教，意味埋め込み，三層分析，SAT，典拠マーカー

\section{はじめに}

\subsection{本稿の位置づけ}

本稿は、拙稿「意味・文体・典拠マーカーの三層地図による仏教文献探索：宗派別参照群・阿弥陀経二訳・親鸞文献の予備的分析」\cite{ueyama2026three-layer-map}を受け、対象を浄土教文献に狭めて再検討する途中報告である。前稿では、漢訳仏典および日本撰述仏教文献を、意味層、文体・語彙層、典拠マーカー層に分けて可視化する枠組みを試みた。そこでは、意味的に近いこと、文体的に近いこと、典拠・引用経路として近いことは同一ではない、という点を中心命題とした。

本稿では、対象をより狭く取り、浄土教、とくに法然と親鸞の関係に集中する。法然の代表書として『選擇本願念佛集』、親鸞の代表書として『顯淨土眞實教行證文類』を置くことは自然である。両者は浄土教の中心問題を共有するため、意味空間上でも大きく重なることが予想される。したがって、本稿で重要なのは「重なるかどうか」ではなく、重なった上で、どの方向にズレが出るかである。

本稿では、この問いに対し、SAT大正新脩大藏經テキストデータベースから取得した漢文本文を用い、法然・親鸞・浄土教祖師文献の意味地図を作成する。さらに、法然側のはみ出し領域と、親鸞側の巻別・祖師文献近接領域を、意味層、文体語彙層、典拠マーカー層の最小構成で読む。

\subsection{研究上の問い}

本稿の問いは次の三点である。

\begin{enumerate}
  \item 法然『選擇本願念佛集』と親鸞『顯淨土眞實教行證文類』は、SAT漢文ベースの意味埋め込み空間でどの程度重なり、どの方向に分かれるか。
  \item 親鸞の広がりは、道綽、善導、源信、曇鸞、天親、龍樹などの浄土教祖師文献とどのような近傍関係を持つか。
  \item 法然側・親鸞側のはみ出し領域を比較したとき、法然は念仏選択の論証へ、親鸞は信巻・化身土巻を中心とする体系的展開へ分かれるのか。
\end{enumerate}

本稿はまだ途中段階であり、ここでの結論は最終的な教義史的主張ではない。意味埋め込みは、研究上の候補領域を見つける探索地図として用いる。本稿の考察は、埋め込み手法とそこから派生する近傍・語群・マーカー集計にもとづく推測に限定する。既存の歴史研究、引用研究、思想史研究との比較は、この方法の信頼性や限界を測るうえで重要であるが、本稿では扱わない。

\section{データ}

\subsection{対象文献}

中心対象は、法然『選擇本願念佛集』と親鸞『顯淨土眞實教行證文類』である。いずれもSAT大正新脩大藏經テキストデータベースから取得した漢文本文を用いた。比較アンカーとして、浄土教祖師文献から龍樹、天親、曇鸞、道綽、善導、源信のテキストを加えた。

\begin{table}[H]
\centering
\small
\caption{対象文献とチャンク数}
\label{tab:corpus}
\begin{tabular}{@{}llllr@{}}
\toprule
分類 & 著者 & 文献 & SAT範囲・ID & chunks \\
\midrule
中心 & 法然 & 選擇本願念佛集 & T2608 & 69 \\
中心 & 親鸞 & 顯淨土眞實教行證文類 & T2646 & 191 \\
アンカー & 龍樹 & 十住毘婆沙論 易行品 & T1521範囲 & 6 \\
アンカー & 天親 & 無量寿経優波提舎願生偈 & T1524 & 9 \\
アンカー & 曇鸞 & 往生論註 & T1819 & 31 \\
アンカー & 道綽 & 安楽集 & T1958 & 42 \\
アンカー & 善導 & 観無量寿仏経疏 & T1753 & 23 \\
アンカー & 源信 & 往生要集 & T2682 & 73 \\
\bottomrule
\end{tabular}
\end{table}

SAT本文の取得元URL、取得条件、処理結果はローカルの出典記録に保存した。本文そのものは公開しない。公開または共有を想定するのは、本文を含まない図、座標、チャンクID、SAT line range、SHA-256 hash、語群カウント、近傍集計に限定する。

\subsection{本文品質と公開境界}

本文取得では、SAT HTMLから本文行を抽出し、UI部品、HTMLタグ、画像プレースホルダ、SAT line referenceの本文混入を除去した。準備済み本文の品質チェックでは、法然T2608が1644 SAT lines、27246 compact chars、親鸞T2646が4215 SAT lines、76030 compact charsであり、\texttt{U+FFFD}、HTMLタグ、HTML entity、画像プレースホルダ、ボタン文字列、SAT line reference混入はいずれも0であった。

前稿で用いた固定長トークン列をそのまま \texttt{encoder.decode()} する方法について、本稿では比較検討を行った。その結果、チャンク境界で \texttt{U+FFFD} が発生する場合がある一方、旧方式とUnicode-safe方式の同一位置チャンクの埋め込み類似度は高く、著者別重心もほぼ一致し、大局的な意味地図への影響は小さいことを確認した。したがって前稿の広域的な配置をただちに否定するものではない。ただし、本稿では本文品質、再現性、SAT line rangeとの対応をより確実にするため、\texttt{tiktoken.decode\_with\_offsets()} を用い、トークン位置をPython文字列のUnicode文字境界へ丸める方法を採用した。700トークン・100トークン重なりという大枠は前稿に合わせつつ、以後の基準処理としてUnicode-safe方式を用いる。

\section{方法}

\subsection{チャンク化と埋め込み}

本文は \texttt{tiktoken} の \texttt{cl100k\_base} でトークン化し、700トークン、100トークン重なりでチャンク化した。チャンク本文はUnicode文字境界に合わせて作成し、各チャンクにSAT line range、チャンク番号、SHA-256 hashを付した。

埋め込みにはOpenAI \texttt{text-embedding-3-large} を用いた。本稿で用いたキャッシュ上の各チャンクベクトルは3072次元である。中心文献である法然・親鸞チャンクでPCAをfitし、この3072次元空間を2次元へ投影した。祖師文献アンカーは、同じ3072次元空間から固定PCA面へ投影した。これにより、法然・親鸞の配置を基準にして、各祖師文献がどの領域へ重なるかを可視化した。PCAの第1軸・第2軸の寄与率はそれぞれ0.068562、0.047405であり、2軸合計は約11.6\%である。したがって、2次元図は3072次元のうち一部を表示する探索用の投影であり、最近傍集計やコサイン類似度は元の3072次元空間で計算したものと併読する。

なお、3072次元から2次元への投影方法はPCAに限られない。UMAP、t-SNE、MDSなどを用いれば、局所近傍やクラスタの見え方が変わる可能性がある。本稿では、軸が線形結合として定義され、同じPCA面へ祖師文献アンカーを後から投影しやすく、再実行時の挙動も比較的安定しているため、まずPCAを基準表示として採用した。別の投影法による頑健性確認は今後の方法検証に回す。

また、二つの文献群が最も重なるように2次元へ投影することも、目的を明示すれば意味を持つ。たとえば法然・親鸞の共通核を見たい場合、両者を分ける方向を弱め、共有される近傍構造を強調する投影を補助図として作ることが考えられる。ただし、そのような投影は差異が見えにくくなるよう設計された図であり、重なったこと自体を独立の証拠として扱うことはできない。したがって本稿では、中立的な基準表示としてPCAを置き、共通核強調図や差異強調図は別の目的付き可視化として扱う。

\subsection{三層分析の最小構成}

本稿では、三層分析を次の最小構成として実装した。

\begin{enumerate}
  \item 意味層: \texttt{text-embedding-3-large} の3072次元コサイン類似度、2次元PCA上の距離、著者別重心、最近傍非自己、最近傍アンカー。
  \item 文体語彙層: 辞書ベースの語群カウント。語群には、念仏・称名、選択・本願、正雑・諸行、廃立・取捨、信・三心、願・回向、名号・真実、往生・浄土、方便・化身土、罪・救済を置いた。
  \item 典拠マーカー層: 経名、論釈名、祖師名、引用導入句の辞書カウント。
\end{enumerate}

これは本格的な引用検出ではない。とくに『教行信証』は大量の引用文から構成されるため、ある語や経名が出たことは、親鸞自身の主張と引用元の主張をただちに区別しない。したがって、本稿の典拠マーカー層は、引用・参照関係そのものではなく、次に精査すべき候補を出す探索的proxyである。

\subsection{巻別・章別ラベル}

親鸞チャンクには、SAT line rangeに基づき、序、教巻、行巻、信巻、証巻、真仏土巻、化身土巻のラベルを付けた。法然チャンクには、序、聖道門・浄土門、正雑二行、本願念仏、三輩・一向専念、三心・信心、付属・証誠・選択総結、末尾・奥書系のラベルを付けた。

ただし、チャンクは100トークンの重なりを持つため、巻境界や章境界をまたぐことがある。ラベルは \texttt{line\_start} に基づく近似であり、厳密な章節単位ではない。

\section{結果}

\subsection{法然・親鸞の意味地図}

図\ref{fig:sat-safe-focus-map}に、SAT漢文ベースの法然・親鸞のみの意味地図を示す。PCAは法然・親鸞チャンクのみでfitし、表示範囲も両者の分布に合わせて調整した。

\begin{figure}[H]
  \centering
  \includegraphics[width=\linewidth]{../figures/sat-safe-honen-shinran-focus-map.png}
  \caption{SAT漢文・Unicode-safe chunkによる法然・親鸞の意味地図。法然・親鸞チャンクでPCAをfitし、法然69チャンク、親鸞191チャンクを表示した。}
  \label{fig:sat-safe-focus-map}
\end{figure}

この図では、法然と親鸞は大きく重なるが、親鸞はより広い分布を示す。法然の分布は比較的まとまり、親鸞は共有核から右方向・下方向を含めて広がる。まずこの二者の関係を基準に置くことで、次に重ねる祖師文献アンカーがどの領域に入るかを読みやすくなる。

\subsection{祖師文献アンカーを重ねた意味地図}

図\ref{fig:sat-safe-anchor-map}に、同じPCA面へ龍樹、天親、曇鸞、道綽、善導、源信を投影した図を示す。PCA fit は図\ref{fig:sat-safe-focus-map}と同じく法然・親鸞チャンクのみであり、祖師文献は後から投影している。

\begin{figure}[H]
  \centering
  \includegraphics[width=\linewidth]{../figures/sat-safe-honen-shinran-high-priest-anchor-map.png}
  \caption{SAT漢文・Unicode-safe chunkによる法然・親鸞・祖師文献アンカー地図。法然・親鸞チャンクでPCAをfitし、龍樹、天親、曇鸞、道綽、善導、源信を同じ面へ投影した。}
  \label{fig:sat-safe-anchor-map}
\end{figure}

祖師文献を重ねると、道綽・善導は共有核に近く、源信・曇鸞は親鸞側の広がりと重なる。図だけではなく高次元空間で最近傍を集計すると、法然と親鸞の重なりはさらに明確になる。

\subsection{2次元方向の後付け解釈}

PCA軸は自動的に教義ラベルを持たない。また、軸の符号は再計算によって反転しうる。したがって、図上の上下左右は、それ自体として固定した意味を持つのではなく、端にあるチャンク群の巻・章、語群、典拠マーカー、上位語から後付けで読む必要がある。ここでは、法然・親鸞focus chunksの各座標について15\%/85\%分位の外側を集計し、表\ref{tab:pca-direction-reading}のように概略ラベルを付けた。

\begin{table}[H]
\centering
\small
\caption{2次元PCA方向の概略ラベル}
\label{tab:pca-direction-reading}
\begin{tabular}{@{}llll@{}}
\toprule
方向 & 主な著者・巻/章 & 多い語群 & 概略ラベル \\
\midrule
右 / PC1+ & 親鸞、信巻・化身土巻 & 罪・救済、廃立・取捨 & 罪救済・化身土側の親鸞展開 \\
左 / PC1- & 法然、正雑二行・一向専念・三心 & 信・三心、往生・浄土、正雑・諸行 & 選択論証・共有核 \\
上 / PC2+ & 親鸞、行巻・真仏土巻 & 願・回向、往生・浄土、念仏・称名 & 阿弥陀・光明・名号・願回向 \\
下 / PC2- & 親鸞、信巻・化身土巻 & 信・三心、罪・救済、方便・化身土 & 信/三心・罪救済・方便/外教 \\
\bottomrule
\end{tabular}
\end{table}

この読みでは、横軸はおおまかに、左の「法然を多く含む選択論証・共有核」から、右の「信巻・化身土巻を中心とする親鸞側の展開」へ向かう。縦軸は、上の「阿弥陀・光明・名号・願回向」方向と、下の「信/三心・罪救済・方便/外教」方向を分けるように見える。ただし、これはこのPCA面での端点集計にもとづく記述であり、高次元空間の全構造や歴史的関係をそのまま表すものではない。

\begin{table}[H]
\centering
\small
\caption{法然チャンクの最近傍非自己}
\label{tab:honen-nearest}
\begin{tabular}{@{}lr@{}}
\toprule
最近傍グループ & chunks \\
\midrule
親鸞 & 50 \\
道綽 & 11 \\
源信 & 4 \\
善導 & 3 \\
曇鸞 & 1 \\
\bottomrule
\end{tabular}
\end{table}

\begin{table}[H]
\centering
\small
\caption{親鸞チャンクの最近傍非自己}
\label{tab:shinran-nearest}
\begin{tabular}{@{}lr@{}}
\toprule
最近傍グループ & chunks \\
\midrule
法然 & 48 \\
道綽 & 45 \\
源信 & 36 \\
曇鸞 & 23 \\
天親 & 21 \\
善導 & 17 \\
龍樹 & 1 \\
\bottomrule
\end{tabular}
\end{table}

表\ref{tab:honen-nearest}では、法然チャンク69件のうち50件が親鸞を最近傍非自己とする。この結果は、少なくとも今回の本文・チャンク・埋め込み条件では、法然と親鸞が意味空間上で強く重なることを示す。ただし、埋め込みの近さそのものは影響関係を証明しない。本稿では影響史の説明へ踏み込まず、意味空間上の近傍関係として扱う。

表\ref{tab:shinran-nearest}では、親鸞チャンク191件の最近傍非自己が法然だけでなく、道綽、源信、曇鸞、天親、善導へ広く分かれる。親鸞は法然との共有核を持つ一方で、祖師文献群の複数方向へ展開している。

\subsection{親鸞の巻別・近接ゾーン}

親鸞チャンクを巻別に分け、最近傍非自己と2次元上の孤立度から、法然近接、祖師近接、親鸞独自候補に分類した。表\ref{tab:shinran-volume-affinity}に主要な巻別・近接ゾーンを示す。

\begin{table}[H]
\centering
\small
\caption{親鸞の巻別・近接ゾーン上位}
\label{tab:shinran-volume-affinity}
\begin{tabular}{@{}llr@{}}
\toprule
巻 & zone & chunks \\
\midrule
信巻 & 法然近接 & 21 \\
化身土巻 & 道綽近接 & 19 \\
行巻 & 法然近接 & 14 \\
信巻 & 親鸞独自候補 & 13 \\
証巻 & 天親近接 & 11 \\
化身土巻 & 法然近接 & 10 \\
化身土巻 & 源信近接 & 9 \\
信巻 & 源信近接 & 8 \\
真仏土巻 & 善導近接 & 8 \\
行巻 & 源信近接 & 8 \\
\bottomrule
\end{tabular}
\end{table}

信巻には法然近接と親鸞独自候補がともに多く現れる。これは、法然から受けた専修念仏の枠組みが、親鸞において信、三心、真実心、回向、本願力の問題へ展開することと対応しうる。化身土巻では道綽近接、源信近接、善導近接、親鸞独自候補が混在する。これは、方便、末法、外教批判、護法・鬼神、真仮の整理などが、単一の祖師文献方向ではなく複合的に広がることを示す。

\subsection{法然三層chunk sequence}

図\ref{fig:honen-three-layer}に、法然『選擇本願念佛集』69チャンクを横軸に取り、意味層、文体語彙層、典拠マーカー層を縦に並べた三層ヒートマップを示す。意味層は、各法然チャンクから親鸞および祖師文献各群への3072次元コサイン類似度の最大値である。文体語彙層と典拠マーカー層は、親鸞図と同じ辞書ベースの語群・明示マーカーを用いた。

\begin{figure}[H]
  \centering
  \includegraphics[width=\linewidth]{../figures/honen-three-layer-sequence-heatmap.png}
  \caption{法然『選択集』三層chunk sequence。上段は意味層、中段は文体語彙層、下段は典拠マーカー層である。横軸は法然チャンク順、縦線はSAT line rangeベースの章・位置境界を示す。本文は含めない。}
  \label{fig:honen-three-layer}
\end{figure}

この図では、意味層の多くが親鸞へ近く出る一方、章・位置によって道綽、善導、源信へ寄る箇所も分かれる。文体語彙層では、正雑二行、本願念仏、三輩・一向専念、三心・信心、付属・証誠・選択総結に、念仏・称名、選択・本願、正雑・諸行、信・三心、往生・浄土の語群が帯状に現れる。典拠マーカー層では引用導入句が全体に強く、七高僧名や経名の出方だけでは章構成を十分に分けられない。したがって法然側の三層図は、法然が親鸞と意味的に重なりつつ、『選択集』内部では念仏を諸行から選び出す論証へ収束していく様子を読むための作業図である。

\subsection{親鸞三層chunk sequence}

図\ref{fig:shinran-three-layer}に、親鸞『教行信証』191チャンクを横軸に取り、意味層、文体語彙層、典拠マーカー層を縦に並べた三層ヒートマップを示す。意味層は、各親鸞チャンクから法然および祖師文献各群への3072次元コサイン類似度の最大値である。文体語彙層は、辞書ベースの語群カウントである。典拠マーカー層は、浄土三部経、涅槃経、大集/日月蔵、論/論註、安楽集、観経疏/善導釈、往生要集、七高僧名、引用導入句の明示マーカーである。

\begin{figure}[H]
  \centering
  \includegraphics[width=\linewidth]{../figures/shinran-three-layer-sequence-heatmap.png}
  \caption{親鸞『教行信証』三層chunk sequence。上段は意味層、中段は文体語彙層、下段は典拠マーカー層である。横軸は親鸞チャンク順、縦線は巻境界を示す。本文は含めない。}
  \label{fig:shinran-three-layer}
\end{figure}

この図では、意味層と文体語彙層が必ずしも同じ動きをしない。たとえば信巻では、意味層の最近傍は法然、源信、道綽に分かれる一方、文体語彙層では信・三心と罪・救済が強く出る。化身土巻では、意味層で道綽、法然、源信が混在し、文体語彙層では信・三心、願・回向、往生・浄土、方便・化身土、廃立・取捨が分散する。典拠マーカー層では引用導入句が全巻を通じて強く、特定典拠名だけでは『教行信証』の引用構造を十分に分けられないことも分かる。したがって、この三層図は、意味的近さ、語彙的特徴、明示的引用マーカーが一致する箇所とズレる箇所を探すための作業図である。

\subsection{法然のはみ出し領域}

法然チャンクのうち、親鸞・祖師文献アンカーから相対的に離れる領域を抽出した。抽出には、2次元PCA上の最近傍非法然距離と、高次元コサイン類似度を組み合わせたprotrusion scoreを用いた。表\ref{tab:honen-protrusion}に上位チャンクを示す。本文は転載せず、SAT line range、章ラベル、最近傍、語群のみを示す。

\begin{table}[H]
\centering
\small
\caption{法然はみ出し領域の上位チャンク}
\label{tab:honen-protrusion}
\begin{tabular}{@{}rlllll@{}}
\toprule
idx & 位置 & line start & 近傍 & anchor & style \\
\midrule
7 & 正雑二行 & T2608\_.83.0003a04 & 親鸞 & 善導 & 正雑・諸行 \\
44 & 付属・証誠・選択総結 & T2608\_.83.0013b17 & 親鸞 & 曇鸞 & 念仏・称名 \\
22 & 三輩・一向専念 & T2608\_.83.0007b07 & 親鸞 & 善導 & 念仏・称名 \\
29 & 三心・信心 & T2608\_.83.0009b05 & 親鸞 & 曇鸞 & 念仏・称名 \\
59 & 付属・証誠・選択総結 & T2608\_.83.0017c18 & 道綽 & 道綽 & 念仏・称名 \\
56 & 付属・証誠・選択総結 & T2608\_.83.0017a01 & 親鸞 & 道綽 & 念仏・称名 \\
15 & 本願念仏 & T2608\_.83.0005b05 & 親鸞 & 源信 & 往生・浄土 \\
37 & 三心・信心 & T2608\_.83.0011b21 & 親鸞 & 源信 & 信・三心 \\
21 & 三輩・一向専念 & T2608\_.83.0007a11 & 道綽 & 道綽 & 正雑・諸行 \\
55 & 付属・証誠・選択総結 & T2608\_.83.0016b22 & 親鸞 & 道綽 & 往生・浄土 \\
\bottomrule
\end{tabular}
\end{table}

法然のはみ出し領域には、正雑二行、三輩・一向専念、本願念仏、三心・信心、付属・証誠・選択総結が多い。これは、念仏そのものの単純な称揚というより、諸行の中から念仏を選び出し、専修の行として位置づける論証の密度を反映していると考えられる。

\subsection{親鸞のはみ出し領域}

同じ尺度で、親鸞チャンクについても非親鸞点からの2次元距離と高次元コサイン類似度を組み合わせ、親鸞側のはみ出し領域を抽出した。表\ref{tab:shinran-protrusion}に上位チャンクを示す。ここでも本文は転載せず、巻、SAT line range、最近傍、語群、内容ゾーンのみを示す。

\begin{table}[H]
\centering
\small
\caption{親鸞はみ出し領域の上位チャンク}
\label{tab:shinran-protrusion}
\begin{tabular}{@{}rlllll@{}}
\toprule
idx & 巻 & zone & line start & 近傍 & style \\
\midrule
162 & 化身土巻 & 護法・鬼神・宇宙秩序 & T2646\_.83.0635b21 & 源信 & 往生・浄土 \\
84 & 信巻 & 信/三心・罪救済 & T2646\_.83.0614a08 & 源信 & 罪・救済 \\
72 & 信巻 & 信/三心・罪救済 & T2646\_.83.0610b22 & 源信 & 罪・救済 \\
83 & 信巻 & 信/三心・罪救済 & T2646\_.83.0613c14 & 道綽 & 罪・救済 \\
85 & 信巻 & 信/三心・罪救済 & T2646\_.83.0614b02 & 源信 & 廃立・取捨 \\
71 & 信巻 & 信/三心・罪救済 & T2646\_.83.0610a27 & 源信 & 願・回向 \\
182 & 化身土巻 & 外教批判 & T2646\_.83.0640b28 & 道綽 & 廃立・取捨 \\
75 & 信巻 & 信/三心・罪救済 & T2646\_.83.0611b10 & 道綽 & 罪・救済 \\
74 & 信巻 & 信/三心・罪救済 & T2646\_.83.0611a10 & 源信 & 罪・救済 \\
78 & 信巻 & 信/三心・罪救済 & T2646\_.83.0612b05 & 源信 & 罪・救済 \\
\bottomrule
\end{tabular}
\end{table}

上位24件を集計すると、巻別には信巻13件、化身土巻8件、真仏土巻2件、証巻1件であった。内容ゾーンでは、信巻の信・三心・罪救済が13件、化身土巻の護法・鬼神・宇宙秩序が4件、真仏土巻の名号・本願・真実が2件、化身土巻の外教批判が1件、方便・真仮整理が1件であった。最近傍非自己は源信10件、道綽6件、曇鸞4件、法然3件、善導1件であり、親鸞のはみ出しは法然から単に離れるというより、源信・道綽・曇鸞などの引用・主題圏と接しながら、信巻と化身土巻の問題群へ広がる。

\subsection{法然・親鸞のはみ出し比較}

法然はみ出し上位24件は、付属・証誠・選択総結10件、三輩・一向専念4件、本願念仏4件、三心・信心3件、正雑二行3件であった。これに対して親鸞は、信巻13件、化身土巻8件を中心に広がる。したがって、同じ「共有核からのはみ出し」でも、法然側は念仏を諸行から選び出す論証の密度として現れ、親鸞側は信巻における罪と救済、化身土巻における方便・真仮・外教批判・護法秩序のような、体系化の周縁部として現れる。

三層sequenceで見ても、この違いはおおむね対応する。法然側では、意味層で親鸞との近さを保ちつつ、文体語彙層が選択・本願、正雑・諸行、念仏・称名、往生・浄土へまとまる。親鸞側では、意味層が法然、源信、道綽などへ分かれ、文体語彙層も信・三心、罪・救済、方便・化身土、廃立・取捨へ広がる。つまり、法然の三層図は論証の収束を、親鸞の三層図は典拠的体系の拡張を示す作業図として読める。

\section{考察}

\subsection{解釈の範囲}

以下の考察は、埋め込み空間上の配置、高次元コサイン類似度、2次元PCA上のはみ出し、辞書ベースの語群、典拠マーカーから得た推測である。歴史的な研究成果、引用研究、教義史的解釈を参照して本稿の読みを補強することは、ここでは行わない。むしろ将来的には、こうした外部研究との照合によって、本手法がどの程度信頼できるか、どこで誤読や過剰解釈を生むかを検証する必要がある。本稿はその前段階として、方法そのものがどのような候補領域を出すかを記述する。

\subsection{法然は「念仏を選ぶ論理」を構成する}

法然『選択集』の特徴は、念仏を諸行から選び出す論証にある。今回のはみ出し領域でも、正雑二行、諸行との取捨、本願、三輩・一向専念、付属・証誠が強く現れた。したがって、法然の独自領域は、いまの段階では「念仏の意味が親鸞と異なる」というより、「なぜ諸行ではなく念仏なのか」を構成する分類・廃立の論証として読むのが自然である。

この読みは、『選択集』という書物の形式とも合う。『選択集』は、浄土三部経や善導などを引きながら、称名念仏を往生行として選び出す書である。親鸞と共有する内容は多いが、書物としての焦点は比較的鋭く、念仏選択の論理へ収束する。

\subsection{親鸞は「本願・信・証の典拠的体系」へ展開する}

親鸞『教行信証』は、今回の意味地図では、法然に近い領域を持ちながら、道綽、源信、曇鸞、天親、善導へ広がった。とくに信巻や化身土巻では、法然近接と親鸞独自候補、祖師近接が混在する。

このことから、親鸞は単に法然を反復するのではなく、法然の念仏選択を、本願、信、回向、名号、真仏土、化身土、方便、罪と救済、外教批判などを含む大きな典拠的体系へ再配置していると読める。とくに親鸞側のはみ出し領域が信巻と化身土巻に集中したことは、この再配置が中心教義の反復だけではなく、救済可能性の限界事例、方便と真実の区別、外教批判や護法秩序まで含む体系拡張として現れていることを示唆する。

\subsection{三層分析の必要性}

ただし、意味層だけでは解釈が混ざる。たとえば、親鸞のあるチャンクが曇鸞に近いとしても、それは曇鸞の引用文が含まれるからなのか、親鸞自身の論述が曇鸞的主題を展開しているからなのか、現段階では区別できない。同様に、法然のはみ出し領域に善導や道綽が近い場合も、実際の引用、語彙、主題、書物構成が混在している。

したがって、本稿の次段階では、意味層、文体語彙層、典拠マーカー層をより明確に分ける必要がある。意味層は主題の近さを示す。文体語彙層は、選択、正雑、廃立、信、願、回向、名号、真実などの語彙操作を示す。典拠マーカー層は、経名、人名、引用導入句、固有句によって、どの文献群を明示的に参照しているかを示す。この三層が一致する箇所とズレる箇所を区別したうえで、必要なら歴史研究との照合へ進む。

\section{限界}

本稿には少なくとも次の限界がある。

\begin{enumerate}
  \item 2次元PCAは全分散の約11.6\%のみを示すため、図上の距離は高次元空間の近傍関係と必ずしも一致しない。
  \item 2次元投影法はPCAに限られない。UMAP、t-SNE、MDSなどで同じ傾向が出るかは未検証である。
  \item チャンクは700トークン・100トークン重なりであり、巻・章・引用単位などの自然境界とは一致しない。
  \item 巻別・章別ラベルはSAT line startに基づく近似であり、境界をまたぐチャンクがある。
  \item 語群カウントは辞書ベースの最小実装であり、否定、引用者、文脈、語義の違いを見ていない。
  \item 典拠マーカー層は引用検出そのものではない。引用文の立場と著者自身の立場を分ける処理が必要である。
  \item 本稿の考察は埋め込み手法等による結果からの推測であり、既存の歴史研究成果によって裏づけたものではない。
  \item 古写経本、異本、写本系資料はまだ直接分析に組み込んでいない。それらの研究成果との照合は、本手法の信頼性や限界を検証する次段階の課題である。
\end{enumerate}

\section{今後の作業}

次に行うべき作業は、第一に、親鸞『教行信証』の巻別・引用別ラベルを精密化することである。特に、引用文と親鸞自身の文を分けることが重要である。第二に、法然『選択集』の章構造を機械的なline rangeではなく、実際の十六章門に対応させる必要がある。第三に、法然側のコーパスを広げる場合は、『黒谷上人語灯録』系、三部経釈系を、本文品質と伝来上の注意を確認した上で扱うべきである。

第四に、公開viewerを作る場合は、本文を出さず、座標、文献ラベル、巻別、近傍表、語群カウント、典拠マーカー集計を表示する形にする。本文はローカル研究用キャッシュとして扱い、公開成果物からは除く。

\section{結論}

本稿では、法然『選択集』、親鸞『教行信証』、および浄土教祖師文献を、SAT漢文・Unicode-safe chunk・\texttt{text-embedding-3-large} によって意味地図化した。法然と親鸞は意味空間上で強く重なるが、親鸞は道綽、源信、曇鸞、天親、善導など複数の祖師文献方向へ広がる。法然のはみ出し領域は、念仏理解の単純な差というより、念仏を諸行から選び出す論証、すなわち正雑二行、取捨、選択、本願、付属・証誠のまとまりとして現れる。親鸞のはみ出し領域は、信巻の罪と救済、化身土巻の方便・真仮整理、護法・鬼神・宇宙秩序、外教批判として現れ、法然側とは異なる広がり方を示した。

したがって、現段階の作業仮説は次のようにまとめられる。法然は念仏を選ぶ論理を構成し、親鸞はその念仏を本願・信・証・真仏土/化身土の典拠的体系として展開する。ただし、これは意味層と語群集計から導いた探索的仮説であり、歴史的研究成果を参照して確認した結論ではない。歴史研究・引用研究との比較は、本手法の信頼性を検証するための別課題として残す。

\begin{thebibliography}{99}

\bibitem{ueyama2026three-layer-map}
上山大信「意味・文体・典拠マーカーの三層地図による仏教文献探索：宗派別参照群・阿弥陀経二訳・親鸞文献の予備的分析」2026.
\url{https://dueyama.github.io/buddhist-text-embedding-map/paper/}

\bibitem{sat}
SAT大正新脩大藏經テキストデータベース.
\url{https://21dzk.l.u-tokyo.ac.jp/SAT/}

\bibitem{jodoshuzensho-senchaku}
新纂浄土宗大辞典「選択本願念仏集」.
\url{https://jodoshuzensho.jp/daijiten/index.php/%E9%81%B8%E6%8A%9E%E6%9C%AC%E9%A1%98%E5%BF%B5%E4%BB%8F%E9%9B%86}

\bibitem{jodoshuzensho-gotoroku}
新纂浄土宗大辞典「黒谷上人語灯録」.
\url{https://jodoshuzensho.jp/daijiten/index.php/%E9%BB%92%E8%B0%B7%E4%B8%8A%E4%BA%BA%E8%AA%9E%E7%81%AF%E9%8C%B2}

\bibitem{jodoshu-senchaku}
浄土宗公式サイト「法然上人のご生涯と浄土宗の教え 選択集」.
\url{https://jodo.or.jp/newspaper/special/6226/}

\bibitem{fujiwara2020}
藤原智「『教行信証』における引用文について：古写経本による再検討」『近現代『教行信証』研究検証プロジェクト研究紀要』3, 63--107, 2020.
\url{https://www.jstage.jst.go.jp/article/kyogyoshinsho/3/0/3_63/_article/-char/ja/}

\bibitem{openai-embeddings}
OpenAI, Embeddings documentation.
\url{https://platform.openai.com/docs/guides/embeddings}

\bibitem{tiktoken}
OpenAI, \texttt{tiktoken}.
\url{https://github.com/openai/tiktoken}

\end{thebibliography}

\end{document}
